\hojaejercicio{Hoja 4}

% Ejercicio 1
\ejercicio{}{}

% Ejercicio 2
\ejercicio{}{}

% Ejercicio 3
\ejercicio{}{}

% Ejercicio 4
\ejercicio{Método de relajación para matrices de diagonal estrictamente dominante.\\
\begin{enumerate}[label=\alph*)]
    \item Probar que si $0 < w \leq 1$ y $\lambda \in \mathbb{C}$ con $|\lambda| \geq 1$ entonces
    \[
        \left| \frac{1 - w - \lambda}{\lambda w} \right| \geq 1.
    \]
    \textit{Indicación: Utilizar que $|x - y| \geq \big||x| - |y|\big|$, $x, y \in \mathbb{C}$.}

    \item Demostrar que si $A$ es de diagonal estrictamente dominante entonces el método de relajación para $A$ es convergente si $0 < w \leq 1$.
\end{enumerate}
}{
    El apartado $a)$ está probado en clase, veamos el apartado $b)$.\\
    Probemoslo por absurdo, es decir, supongamos que:
    \[
        \exists \lambda \in \sigma(L_R) \text{ tal que } |\lambda| \geq 1 \implies
        \rho(L_R) \geq 1.
    \]
    Ahora entonces:
    \[
        \det(L_R - \lambda I) =
        0 =
        \det( \left( \frac{D}{\omega} - E \right)^{-1} \left( F + \left( \frac{1 - \omega}{\omega} \right) D \right) - \lambda I ) 
    \]
}

% Ejercicio 5
\ejercicio{Teorema de los círculos de Gershgorin y matrices de diagonal estrictamente dominante.\\
\begin{enumerate}[label=\alph*)]
    \item \textbf{Teorema de los círculos de Gershgorin:} Si $A = (a_{ij})_{i,j=1}^n \in \mathcal{M}_n$, demostrar que
    \[
        \mathrm{sp}(A) \subset \bigcup_{i=1}^n \overline{B}_{r_i}(a_{ii})
        \quad \text{siendo} \quad
        r_i = \sum_{\substack{j=1\\j \neq i}}^n |a_{ij}|,
    \]
    donde
    \[
        \overline{B}_{r_i}(a_{ii}) = \{z \in \mathbb{C} : |z - a_{ii}| \leq r_i\}
    \]
    es la bola cerrada de centro $a_{ii}$ y radio $r_i$.
    
    \textit{Indicación: Téngase en cuenta que las matrices de diagonal estrictamente dominante son inversibles.}

    \item Sea $A = (a_{ij})_{i,j=1}^n \in \mathcal{M}_n$ una matriz de diagonal estrictamente dominante con $a_{ii} > 0$, $i = 1, \ldots, n$ y $\mathrm{sp}(A) \subset \mathbb{R}$.
    
    Utilizar el apartado a) para probar que sus autovalores son estrictamente positivos.

    \item Sea $A = (a_{ij})_{i,j=1}^n \in \mathcal{M}_n$ una matriz hermítica, de diagonal estrictamente dominante con $a_{ii} > 0$, $i = 1, \ldots, n$.
    
    Probar que el método de relajación para $A$ es convergente si, y sólo si, $0 < \omega < 2$.
\end{enumerate}
}{
    Veamos los diferentes apartados.\\
    \begin{enumerate}[label=\alph*)]
        \item Si $\det(A - \lambda I) = 0$ entonces $A - \lambda I$ no es d.e.d., por lo que existe $k \in \{1, \ldots, n\}$ tal que:
        \[
            |b_{kk}| \leq \sum_{\substack{j=1\\j \neq k}}^n |b_{kj}| \implies
            |\lambda - a_{kk}| \leq \sum_{\substack{j=1\\j \neq k}}^n |a_{kj}| =
            r_k. \implies
        \]
        \[
            \lambda \in \overline{B}_{r_k}(a_{kk}) \subset \bigcup_{i=1}^n \overline{B}_{r_i}(a_{ii}).
        \]
        \item Dado que $a_{ii} > 0$ y $r_i < a_{ii}$, al ser una matriz d.e.d., entonces $a_{ii} - r_i > 0$ y por lo tanto $\lambda \in \overline{B}_{r_i}(a_{ii})$ implica que $\lambda > 0$.
        \item Visto en teoría.
    \end{enumerate}
}